\section{Purpose}
The purpose of this document (Implementation and Test Deliverable - ITD) is to describe the implementation of the Best Bike Paths (BBP) software system, specifying how the requirements defined in the RASD and the architecture designed in the DD have been translated into source code.

This document details the functionalities implemented in the final prototype, as well as the technologies and frameworks used for development. Furthermore, it describes the code structure, the strategies adopted to verify its correctness, and the necessary instructions for installation and execution.

In particular, this document highlights deviations from the original design, justifying the implementation choices made to adapt the project to time and technical constraints. It also describes the tools specifically developed to demonstrate the system's functionality.

\section{Scope}
The implemented system constitutes a functional prototype of the Best Bike Paths mobile application.
The main objective of this implementation is to provide users (cyclists and tourists) with a tool to track their cycling activities, manage their route history, and visualize interactive maps.

Compared to the complete vision described in the RASD, the scope of this release focuses on the "Core" functionalities:
\begin{itemize}
    \item \textbf{User Management:} Registration, authentication, and profile management.
    \item \textbf{Activity Tracking:} Real-time GPS tracking, calculation of basic statistics (distance, time, speed), and map visualization.
    \item \textbf{Content Management:} Route saving and history visualization.
    \item \textbf{Route Search:} Searching for routes visible to the user by entering departure and destination points.
\end{itemize}

\paragraph{Exclusions and Limitations}
In order to ensure the stability and quality of the application's main core, and in accordance with the Assignment requirements, some advanced features described in the design have been excluded from this version of the prototype:
\begin{itemize}
    \item \textbf{Report Merge Logic:} The complex algorithm for the automatic unification of user reports has not been implemented.
    \item \textbf{Automatic Detection (Sensor-based):} Automatic pothole detection via accelerometer is not included.
    \item \textbf{Weather Integration:} External meteorological data is not retrieved.
\end{itemize}

\paragraph{Validation Support Features}
To allow for effective system verification in static test environments (e.g., classroom presentation) where physical GPS movement is not possible, the project scope has been extended to include a Route Simulator (Mocking System). This feature allows the application to generate fictitious GPS tracks and simulate user movement, enabling a complete demonstration of the "Start Activity $\rightarrow$ Tracking $\rightarrow$ End Activity $\rightarrow$ Data Analysis" flow without the need for physical displacement.

\section{Definitions, Acronyms, and Abbreviations}
\subsection{Definitions}
\begin{itemize}
    \item \textbf{Backend-as-a-Service (BaaS)}: A cloud computing model adopted in this architecture where backend responsibilities (database, authentication, hosting) are outsourced to a managed provider (Supabase), allowing development to focus on frontend logic.
    \item \textbf{Cron Job}: A time-based job scheduler in Unix-like computer operating systems, used in this project to trigger periodic maintenance tasks (e.g., the decay of old feedbacks) on the server.
    \item \textbf{Edge Function}: A piece of serverless application logic (e.g., the path merging algorithm) that runs on the cloud provider's infrastructure, triggered by specific database events or HTTP requests.
    \item \textbf{Hybrid Application}: The chosen architectural approach where the mobile app is built using web technologies (React/TypeScript) but executes within a native container to access device hardware.
    \item \textbf{Kalman Filter}: An algorithm that uses a series of measurements observed over time (containing statistical noise) to produce estimates of unknown variables (e.g., precise GPS coordinates) that tend to be more accurate than those based on a single measurement alone.
    \item \textbf{Native Bridge}: The software mechanism (provided by Capacitor) that allows the web-based client code to communicate with the native operating system and access sensors like the Accelerometer and GPS in the background.
    \item \textbf{Object Storage}: A computer data storage architecture that manages data as distinct units (objects), used in this system to store large binary assets like user profile pictures.
    \item \textbf{PostGIS}: An open-source software extender for the PostgreSQL database, used in this project to store, index, and query complex geospatial data (geometries) efficiently.
    \item \textbf{Row Level Security (RLS)}: A security mechanism implemented directly in the database engine that restricts access to specific data rows based on the user's authentication token, eliminating the need for application-level security filters.
    \item \textbf{Thick Client}: An architectural style where the client application (running on the user's device) handles the bulk of data processing, validation, and temporary storage, contrasting with a Thin Client that relies heavily on the server.
    \item \textbf{Webhook}: A mechanism for an application to provide other applications with real-time information. In this system, database webhooks trigger Edge Functions automatically when data is inserted or updated.
\end{itemize}

\subsection{Acronyms}
\begin{itemize}
    \item \textbf{API}: Application Programming Interface.
    \item \textbf{BaaS}: Backend-as-a-Service.
    \item \textbf{CPU}: Central Processing Unit.
    \item \textbf{CRUD}: Create, Read, Update, Delete.
    \item \textbf{JSON}: JavaScript Object Notation.
    \item \textbf{JWT}: JSON Web Token (used for stateless authentication).
    \item \textbf{KPI}: Key Performance Indicator.
    \item \textbf{OS}: Operating System.
    \item \textbf{PWA}: Progressive Web Application.
    \item \textbf{RLS}: Row Level Security.
    \item \textbf{RPC}: Remote Procedure Call (used to invoke Edge Functions).
    \item \textbf{SDK}: Software Development Kit.
    \item \textbf{SPA}: Single Page Application.
    \item \textbf{SQL}: Structured Query Language.
    \item \textbf{WSS}: WebSocket Secure (used for real-time updates).
\end{itemize}

\subsection{Abbreviations}
\begin{itemize}
    \item \textbf{DB}: DataBase.
    \item \textbf{UI}: User Interface.
    \item \textbf{VM}: Virtual Machine.
\end{itemize}

\section{Revision History}
Version 1.0: 01/02/2026

\section{Reference Documents}
\begin{itemize}
    \item Requirement Analysis and Specification Document (RASD) v2.0
    \item Design Document (DD) v1.0
    \item Assignment RDD AY 2025-2026 - Software Engineering 2 Project Assignment, Politecnico di Milano.
\end{itemize}

\section{Document Structure}
This document is organized into the following chapters:

\begin{itemize} 
    \item \textbf{Chapter 1 Introduction:}
    Provides an overview of the document, defining the purpose and scope of the implementation. It explicitly states the boundaries of the prototype, highlighting the focus on core functionalities and the introduction of a "Demo Mode" for validation purposes in static environments.

    \item \textbf{Chapter 2 Implemented Requirements and Functions:}
    Details the functional requirements that have been successfully implemented (e.g., User Management, Activity Tracking) and lists the functionalities excluded from this release (such as automatic sensor detection and merge logic), providing motivations based on project constraints. It also describes the additional "Mocking System" developed to support the demonstration.

    \item \textbf{Chapter 3 Adopted Development Frameworks and Technologies:}
    Describes the technology stack used to build the solution (including React, TypeScript, Capacitor, and Supabase), referencing the choices made in the Design Document and explaining the advantages of the adopted middleware and tools.

    \item \textbf{Chapter 4 Source Code Structure:} 
    Explains the organization of the source code repository, detailing the folder hierarchy and the location of the main architectural components.

    \item \textbf{Chapter 5 Testing:} 
    Outlines the testing strategies employed to ensure software quality, describing the procedures used for system testing and the validation of the core user flows through the implemented prototype.

    \item \textbf{Chapter 6 Installation Instructions:} 
    Provides step-by-step instructions for setting up the development environment, installing necessary dependencies, and executing the software to reproduce the implemented features.

    \item \textbf{Chapter 7 Effort Spent:} 
    Summarizes the time effort dedicated by each team member to the development of the code and the writing of this document.
\end{itemize}

\section{Use of Generative AI Tools}
In compliance with the course guidelines regarding academic integrity and content responsibility, we declare that we have used Generative Artificial Intelligence (GenAI) tools to support the drafting of this document (Design Document - DD).

We emphasize that the entire content of this document is the result of our original reasoning; AI was not used to generate architectural decisions or core logic from scratch, but exclusively as a tool for linguistic support, formal verification, and visual prototyping. The details of its usage are reported below:

\begin{enumerate}
    \item \textbf{Tools Used}
    \begin{itemize}
        \item \textbf{Textual Support:} ChatGPT (model GPT-3.5/4o mini), Gemini Advanced (Gemini 3).
        \item \textbf{Visual Generation (UI Mockups):} Lovable.
    \end{itemize}

    \item \textbf{Provided Inputs (Prompts and Context)} The inputs provided to the system were exclusively textual and based on preliminary work carried out by the group. Specifically:
    \begin{itemize}
        \item \textbf{For linguistic revision:} We input paragraphs or entire sections written by us in Italian/English, with the request to improve grammatical form, academic vocabulary, and text fluidity, without altering the technical meaning.
        \begin{quote}
            \textbf{Example A - For stylistic improvement:} \textit{"Act as a professional Technical Writer reviewing a Software Design Document. I will provide you with a raw text describing the 'Path Merging Algorithm'. Please rewrite it to improve grammatical correctness and use formal academic English suitable for a system architecture description. Do NOT add new information or change the logical flow. Just polish the language."}
        \end{quote}
        
        \item \textbf{For completeness verification:} We provided lists of component interfaces or sequence diagram steps developed by the group, asking the model to flag potential logical ambiguities or missing edge cases.
        \begin{quote}
            \textbf{Example B - For completeness verification:} \textit{"Act as a Senior Software Engineer. Below is the list of APIs we defined for the 'Sensor Controller'. Analyze them for consistency with the REST standard. Specifically, identify if there are any missing response codes or parameters that a standard Design Document should include."}
        \end{quote}
        
        \item \textbf{For UI Prototyping (via Lovable):} We provided detailed descriptions of the application screens, including layout structure, color palettes, and required UI elements, to generate high-fidelity visual mockups for Chapter 3.
        \begin{quote}
            \textbf{Example C - For UI Generation:} \textit{"Create a mobile screen for a cycling app in a modern style. It must include a map view showing a route, a bottom sheet showing real-time speed (25 km/h) and distance (12.4 km), and a large red circular button to 'Report Hazard'."}
        \end{quote}
    \end{itemize}

    \item \textbf{Obtained Outputs}
    \begin{itemize}
        \item Suggestions for sentence reformulation to make the text more formal and concise.
        \item Checklists to verify consistency between different diagrams and textual descriptions.
        \item Graphical assets and UI Mockups representing the application interface, which were subsequently refined to match the UX flow.
    \end{itemize}

    \item \textbf{Verification, Refinement, and Integration} Every output generated by the AI was subjected to a manual critical review by the group members.
    \begin{itemize}
        \item We did not passively accept the proposed reformulations; instead, we selected and adapted them to ensure they exactly reflected our project specifications.
        \item The UI mockups generated by Lovable were manually reviewed and used as a basis to define the final interface design, ensuring consistency with the functional requirements.
    \end{itemize}
\end{enumerate}

